% Part: first-order-logic
% Chapter: axiomatic-deduction
% Section: proving-things

\documentclass[../../../include/open-logic-section]{subfiles}

\begin{document}

\iftag{FOL}
      {\olfileid{fol}{axd}{pro}}
      {\olfileid{pl}{axd}{pro}}

\olsection{د \usetoken{P}{derivation} بېلګې}


\begin{ex}
فرض کړئ غواړو $(\lnot !D \lor !E) \lif (!D \lif !E)$ ثابت کړو۔
دا په څرګند ډول زموږ د هېڅ بديهي اصل بېلګه نۀ ده، نو د هغې د
!!{derive} کولو دپاره د~\MP{} قاعده کارول پکار دي۔ زموږ يوازينۍ
قاعده مودوس پوننس ده، چې $!A$ او $!A \lif !B$ راکړل شي، د~$!B$
توجيه راکوي۔ يوه تګلاره دا ده چې \olref[prp]{ax:lor3} داسې وکاروو
چې $!A$، $\lnot !D$؛ $!B$، $!E$؛ او $!C$، $!D \lif !E$ وي؛ يعنې
دا بېلګه:
\[
(\lnot !D \lif (!D \lif !E)) \lif ((!E \lif (!D \lif !E)) \lif ((\lnot
!D \lor !E) \lif (!D \lif !E))).
\]
ولې؟ د مودوس پوننس دوه تطبيقونه وروستۍ برخه راکوي، او همدا مو
غوښتې ده۔ په اسانۍ وينو چې $\lnot !D \lif (!D \lif !E)$ د
\olref[prp]{ax:lnot2}، او $!E \lif (!D \lif !E)$ د
\olref[prp]{ax:lif1} بېلګه ده۔ نو زموږ !!{derivation} دا دے:
\begin{derivation}
  1. &  $\lnot !D \lif (!D \lif !E)$ & \olref[prp]{ax:lnot2} \\
  2. & $(\lnot !D \lif (!D \lif !E)) \lif {}$\\
  &\qquad $((!E \lif (!D \lif !E)) \lif ((\lnot !D \lor !E) \lif (!D \lif !E)))$ & \olref[prp]{ax:lor3}\\
  3. & $(!E \lif (!D \lif !E)) \lif ((\lnot !D \lor !E) \lif (!D \lif !E))$ & 1, 2, \MP\\
  4. & $!E \lif (!D \lif !E)$ & \olref[prp]{ax:lif1}\\
  5. & $(\lnot !D \lor !E) \lif (!D \lif !E)$ & 3, 4, \MP
\end{derivation}
\end{ex}

\begin{ex}\ollabel{ex:identity}
راځئ د~$!D \lif !D$ يو !!a{derivation} ولټوو۔ دا د بديهي اصل بېلګه
نۀ ده، نو په~\MP{} ئې !!{derive} کول پکار دي۔
\olref[prp]{ax:lif1} د~$!A \lif !B$ په بڼه يو بديهي اصل دے چې
\MP{} پرې پلي کېدے شي۔ د ګټې دپاره، البته، هغه~$!B$ چې \MP{} ئې په
دې حالت کښې د سم ګام په توګه توجيه کوي بايد~$!D \lif !D$ وي؛ ځکه
همدا مو غوښتل چې !!{derive} ئې کړو۔ په دې سره $!A$ هم بايد $!D$ وي؛
يعنې د \olref[prp]{ax:lif1} دې بېلګې ته کتلے شو:
\[
!D \lif (!D \lif !D)
\]
د~\MP{} د پلي کولو دپاره به د دويمې اړوندې مقدمې، يعنې~$!A$،
توجيه هم پکار وي۔ خو زموږ په حالت کښې هغه~$!D$ دے، او $!D$ په
يوازې ځان !!{derive} کولے نۀ شو۔ نو بله تګلاره پکار ده۔

بل بديهي اصل چې يوازې~$\lif$ لري \olref[prp]{ax:lif2} دے، يعنې
\[
(!A \lif (!B \lif !C)) \lif ((!A \lif !B) \lif (!A \lif !C))
\]
د~\MP{} په دوه ځله پلي کولو وروستي دننني شرطي ته رسېدے شو۔ بيا دا
معنٰی لري چې د \olref[prp]{ax:lif2} داسې بېلګه غواړو چې
$!A \lif !C$ پکې $!D \lif !D$، يعنې زموږ موخه !!{formula} وي۔
نو $!A$ او $!C$ دواړه~$!D$ دي۔ $!B$ څنګه وټاکو چې دواړه
$!A \lif (!B \lif !C)$ او $!A \lif !B$، يعنې زموږ په حالت کښې
$!D \lif (!B \lif !D)$ او $!D \lif !B$، هم !!{derivable} وي؟
له دغو لومړنے، هر~$!B$ چې وټاکو، لا د \olref[prp]{ax:lif1} بېلګه
ده۔ او که $!B$، $(!D \lif !D)$ وي، نو $!D \lif !B$ به هم د
\olref[prp]{ax:lif1} بله بېلګه وي۔ نو زموږ !!{derivation} دا دے:
\begin{derivation}
  1. & $!D \lif ((!D \lif !D) \lif !D)$ & \olref[prp]{ax:lif1}\\
  2. & $(!D \lif ((!D \lif !D) \lif !D)) \lif {}$\\
  & \qquad $((!D \lif (!D \lif !D)) \lif (!D \lif !D))$ & \olref[prp]{ax:lif2}\\
  3. & $(!D \lif (!D \lif !D)) \lif (!D \lif !D)$ & 1, 2, \MP\\
  4. & $!D \lif (!D \lif !D)$ & \olref[prp]{ax:lif1}\\
  5. & $!D \lif !D$ & 3, 4, \MP
\end{derivation}
\end{ex}

\begin{ex}\ollabel{ex:chain}
کله کله غواړو وښيو چې له ځينو نورو !!{formula}s~$\Gamma$ څخه د کوم
!!{formula} يو !!a{derivation} شته۔ مثلاً راځئ وښيو چې له
$\Gamma = \{!A \lif !B, !B \lif !C\}$ څخه
$!A \lif !C$، !!{derive} کولے شو۔
\begin{derivation}
  1. & $!A \lif !B$ & \Hyp\\
  2. & $!B \lif !C$ & \Hyp\\
  3. & $(!B \lif !C) \lif (!A \lif (!B \lif !C))$ & \olref[prp]{ax:lif1} \\
  4. & $!A \lif (!B \lif !C)$ & 2, 3, \MP\\
  5. & $(!A \lif (!B \lif !C)) \lif {}$\\
  & \qquad $((!A \lif !B) \lif (!A \lif !C))$ & \olref[prp]{ax:lif2}\\
  6. & $((!A \lif !B) \lif (!A \lif !C))$ & 4, 5, \MP\\
  7. & $!A \lif !C$ & 1, 6, \MP
\end{derivation}
هغه کرښې چې ``\Hyp'' (د ``فرضیې'' دپاره) پرې ليکل شوي دي، ښيي چې
په هغې کرښه !!{formula} د~$\Gamma$ يو !!a{element} دے۔
\end{ex}

\begin{prop}
  \ollabel{prop:chain} که $\Gamma \Proves !A \lif !B$ او $\Gamma
  \Proves !B \lif !C$ وي، نو $\Gamma \Proves !A \lif !C$
\end{prop}

\begin{proof}
  فرض کړئ $\Gamma \Proves !A \lif !B$ او $\Gamma \Proves !B \lif
  !C$۔ بيا له~$\Gamma$ څخه د~$!A \lif !B$ يو !!a{derivation} شته؛
  او له~$\Gamma$ څخه د~$!B \lif !C$ يو !!a{derivation} هم شته۔ د
  هغو په يو بل پسې نښلولو ئې په يو !!{derivation} کښې يو ځاے کړئ۔
  اوس د مخکښې بېلګې د~!!{derivation} ۳--۷ کرښې ورزياتې کړئ۔ دا له~$\Gamma$ څخه د
  $!A \lif !C$ يو !!a{derivation} دے---او همدا د نوي
  !!{derivation} وروستۍ کرښه ده۔ پام وکړئ چې د ۴مې او ۷مې کرښې
  توجيه لا هم سمه پاتې کېږي که د ۱مې کرښې حواله د~$!A \lif !B$
  د~!!{derivation} د وروستۍ کرښې په حواله، او د ۲مې کرښې حواله د
  $!B \lif !C$ د~!!{derivation} د وروستۍ کرښې په حواله بدله شي۔
\end{proof}

\begin{prob} 
له بديهي اصولو څخه د !!{derivation}s په وړاندې کولو وښيئ چې لاندې
خبرې سمې دي:
\begin{enumerate} 
\item $(!A \land !B) \lif (!B \land !A)$
\item $((!A \land !B) \lif !C) \lif (!A \lif (!B \lif !C))$
\item $\lnot(!A \lor !B) \lif \lnot !A$
\end{enumerate} 
\end{prob}



\end{document}
